\section{Nghiên cứu chương sách giáo trình}

\subsection{Thông tin chương sách được chọn}
\begin{itemize}
    \item \textbf{Tên sách:} \textit{[Handbook of Biometrics / Handbook of Face Recognition 2nd Edition]}.
    \item \textbf{Tên chương:} \textit{[Điền tên chương]}.
    \item \textbf{Chủ đề chính:} \textit{[Điền kỹ thuật hoặc bài toán nhận dạng trọng tâm]}.
    \item \textbf{Lý do lựa chọn:} \textit{[Giải thích vì sao chương này phù hợp với đồ án]}.
\end{itemize}

\subsection{Bài toán và giả định trong chương sách}
Trình bày bài toán mà chương sách giải quyết, bao gồm dữ liệu đầu vào, đầu ra mong muốn, giả định về dữ liệu, môi trường triển khai và các ràng buộc kỹ thuật. Phần này cần làm rõ chương sách tập trung vào nhận dạng, xác minh, phát hiện hay phân loại.

\subsection{Kỹ thuật hoặc phương pháp cốt lõi}
Mô tả chính xác phương pháp trọng tâm trong chương sách. Nội dung nên bao gồm:
\begin{itemize}
    \item Ý tưởng trực giác của phương pháp.
    \item Biểu diễn toán học hoặc mô hình kiến trúc nếu có.
    \item Các bước xử lý chính từ dữ liệu đầu vào đến kết quả đầu ra.
    \item Vai trò của từng thành phần trong toàn bộ hệ thống nhận dạng.
\end{itemize}

\subsection{Quy trình hoạt động chi tiết}
Trình bày quy trình theo từng bước. Có thể dùng sơ đồ khối, thuật toán giả hoặc bảng tóm tắt để minh họa luồng xử lý.

\subsubsection{Tiền xử lý dữ liệu}
Mô tả các bước chuẩn hóa, căn chỉnh, lọc nhiễu, phát hiện vùng quan tâm hoặc tăng cường dữ liệu được nêu trong chương sách.

\subsubsection{Trích xuất đặc trưng}
Giải thích cách phương pháp tạo ra đặc trưng phân biệt cho đối tượng cần nhận dạng. Nếu chương sách sử dụng descriptor thủ công hoặc mô hình học máy truyền thống, cần nêu rõ ưu và nhược điểm.

\subsubsection{So khớp, phân loại hoặc xác minh}
Trình bày cách hệ thống đưa ra quyết định cuối cùng, ví dụ dùng khoảng cách, ngưỡng, bộ phân loại hoặc mô hình xác suất.

\subsection{Ưu điểm và hạn chế từ góc nhìn giáo trình}
Tóm tắt các điểm mạnh, điểm yếu và điều kiện áp dụng của phương pháp trong chương sách. Đây là cơ sở để so sánh với bài báo SOTA ở phần sau.